% SPDX-License-Identifier: CC-BY-SA-4.0
% Author: Matthieu Perrin

\begingroup

\SetKwData{Acks}{acks}
\SetKwData{Delivered}{delivered}
\SetKwBroadcast{MB}{mb}
\SetKwBroadcast{RBF}{rb-fifo}
\SetKwMessage{Init}{init}
\SetKwMessage{Ack}{ack}
\SetKwVariable{Payload}{payload}
\SetKwVariable{SN}{sn}
\SetKwVariable{Date}{date}

\begin{frame}{Implémentation de Mutual Broadcast}

  \on[x=-3mm, y=-3mm]{
    \begin{algorithm}[H]
      \LVariables{\At $p_i$}{
        $\Delivered_i \leftarrow [0, ..., 0]$\tcp*[l]{Version vector}
        $\Acks_i \leftarrow 0$\tcp*[l]{Compte pour message courant}
      }
      \Method{$\MB.\Broadcast(\Payload)$ \At $p_i$}{
        \nl $\Acks_i \leftarrow 0$\\
        \nl \Let $\SN \leftarrow \Delivered_i[i]$\\
        \nl \RBF.\Broadcast \alert{$\Init(\Payload, \SN)$}\\
        \nl \Wait $\Delivered_i[i] > \SN$
      }
      \Upon{\RBF.\Deliver $\alert{\Init(\Payload, \SN)}$ \From $p_i$ \At $p_j$}{
        \nl \Send \example{$\Ack(\SN, \Exampleb<3>{\Delivered_j})$} \To $p_i$\\
        \nl \Alertb<2>{\lIf{$j = i$}{
          \Wait $\Acks_j > \frac{n}{2}$
        }}\\
        \nl \Trigger \Alertb<2>{\MB.\Deliver $\Payload$} \From $p_i$\\
        \nl $\Delivered_j[i] \leftarrow \Delivered_j[i] + 1$
      }
      \Upon{\Receive $\example{\Ack(\SN, \Exampleb<3>{\Date})}$ \From $p_j$ \At $p_i$}{
        \nl \Exampleb<3>{\Wait{$\forall k, \Delivered_i[k] \ge \Date[k]$}}\\
        \nl \lIf{$\SN = \Delivered_i[i]$}{$\Acks_i \leftarrow \Acks_i + 1$}
      }
    \end{algorithm}
  }
  
  \on[right=.4\textwidth, y=18mm]{
    \begin{tikzpicture}[y=7mm, x=4mm]

      \draw[process] (0,4) node[left]{$p_1$} to (10,4);
      \draw[process] (0,3) node[left]{$p_2$} to (10,3);
      \draw[process] (0,2) node[left]{$p_3$} to (10,2);
      \draw[process] (0,1) node[left]{$p_4$} to (10,1);
      \draw[process] (0,0) node[left]{$p_5$} to (10,0);
      
      \coordinate    (init1) at (1,4);
      \coordinate    (init2) at (1,0);

      \uncoverb<1>{
        \footnotesize
        \draw[message, alert]     (init1.center) to[bend right]                                    (2.0,4);
        \draw[message, alert]     (init1.center) to                                                (2.6,3);
        \draw[message, alert]     (init1.center) to                                                (3.4,2);
        \draw[message, alert]     (init1.center) to                                                (4.2,1);
        \draw[message, alert]     (init1.center) to             node[sloped, swap] {\Init} (5.0,0);
        
        \draw[message, example] (2.2,4)        to[bend right] node[sloped, swap] {\Ack}  (3.2,4);
        \draw[message, example] (2.8,3)        to             node[sloped, swap] {\Ack}  (3.8,4);
        \draw[message, example] (3.6,2)        to             node[sloped, swap] {\Ack}  (5.6,4);
        \draw[message, example] (4.4,1)        to             node[sloped, swap] {\Ack}  (7.4,4);
        \draw[message, example] (5.2,0)        to             node[sloped, swap] {\Ack}  (9.2,4);
      }

      \uncoverb<2>{
        \node[alert,   operation, fit={(init1.center)(6,4)}] {};

        \draw[message, fade]      (init1.center) to[bend right]                                    (2.0,4);
        \draw[message, fade]      (2.2,4)        to[bend right]                                    (3.2,4);
        \draw[message, fade]      (2.8,3)        to                                                (3.8,4);
        \draw[message, fade]      (3.6,2)        to                                                (5.6,4);
        \draw[message, fade]      (4.4,1)        to                                                (7.4,4);
        \draw[message, fade]      (5.2,0)        to                                                (9.2,4);

        \draw[message, alert]     (init1.center) to[bend left]                                     (6,4);
        \draw[message, alert]     (init1.center) to                                                (2.6,3);
        \draw[message, alert]     (init1.center) to                                                (3.4,2);
        \draw[message, alert]     (init1.center) to                                                (4.2,1);
        \draw[message, alert]     (init1.center) to                                                (5.0,0);
      }
      
      \uncover<3>{
        \node[alert,   operation, fit={(init1.center)(9,4)}] {};
        \node[structure, operation, fit={(init2.center)(3.5,0)}] {};

        \footnotesize
        \draw[message, alert]     (init1.center) to[bend right]                                    (2.0,4);
        \draw[message, alert]     (init1.center) to                                                (2.6,3);
        \draw[message, alert]     (init1.center) to                                                (3.4,2);
        \draw[message, example]   (3.6,2)        to             node[sloped] {[0,0,0,0,1]}         (5.6,4) to[bend right] (9,4);
        \draw[message, example]   (2.2,4)        to[bend right]                                    (3.2,4);
        \draw[message, example]   (2.8,3)        to                                                (3.8,4);

        \draw[message, structure] (init2.center) to[bend right]                                    (2.0,0);
        \draw[message, structure] (init2.center) to                                                (2.0,1);
        \draw[message, structure] (init2.center) to                                                (2.0,2);
        \draw[message, structure] (init2.center) to                                                (8.0,4);
        \draw[message, fade]      (2.2,1)        to                                                (3.0,0);
        \draw[message, fade]      (2.2,2)        to                                                (3.5,0);

        \path[example, thick, ->, shorten <=3pt, shorten >=3pt] (2.2,2) edge[bend right] (3.6,2);
      }
    \end{tikzpicture}
  }

  \on<3>[right=.4\textwidth, anchor=north, y=-9mm]{
    \begin{itemize}
    \item $p_3 \in Q_1 \cap Q_5$
    \item $p_3$ délivre $m_5$ avant $m_1$
    \item $\example{\Ack}$ transmet la causalité
    \item $p_1$ délivre $m_5$ avant $m_1$
    \item pas MP1
    \end{itemize}
  }
  
\end{frame}

\endgroup
\endinput
